\begin{cactuspart}{4}{Appendices}{$RCSfile$}{$Revision$}
\renewcommand{\thepage}{\Alph{part}\arabic{page}}

\chapter{Configuration file syntax}
\label{sec:cofisy}    

\section{General Concepts}
\label{sec:geco}

Each thorn is configured by three compulsory files in the top level thorn
directory:
\begin{itemize}
\item{} interface.ccl
\item{} param.ccl
\item{} schedule.ccl
\end{itemize}
These files are written in the {\it Cactus Configuration Language} which is
case insensitive.

\section{interface.ccl}
\label{sec:in}

The interface configuration file consists of a header block giving details 
of the thorns relationship with other thorns, and then a series of blocks 
listing the thorn's global variables.
The header block has the form:
{\t 
\begin{verbatim}
implements: <implementation>
[inherits: <implementation>, <implementation>]
[friend: <implementation>, <implementation>]
\end{verbatim}
}
where
\begin{itemize}
\item{} The implementation name must be unique among all thorns, except between 
        thorns which have the same public and protected variables and 
        parameters.
\item{} Inheriting from another implementation makes all that implementation 
        public variables available to your thorn. At least one thorn 
        providing any inherited implementation must be present at compile time.
        A thorn cannot inherit from itself. Inheritance is associative and 
        recursive, but not commutative.
\item{} Being a friend of another implementation makes all that 
        implementations 
        protected variables available to your thorn. At least one thorn
        providing an implementation for each friend must be present at
        compile time. A thorn cannot be its own friend. Friendship is 
        associative and commutative. Note that your thorn is also a friend
        of all your thorns friend's friends.
\end{itemize} 

The thorn's variables are defined by:
{\t
\begin{verbatim}
[<access>:]

VTYPE <group name> [TYPE=<gtype>] [DIM=<dim>] 
[{
   <variable name>[,]<variable name>
   <variable name>
}] [``<group description>'']
\end{verbatim}}

\begin{itemize}
\item{} {\t access} defines which thorns have can use the following 
        groups, and can be either {\t public}, {\t protected} or {\t private}.
\item{} {\t VTYPE} defines the data type of the variables in the group.
Supported data types are {\t LOGICAL}, {\t INTEGER}, {\t CHAR} and {\t REAL}.
(In the future {\t COMPLEX} will also be supported.)
\item{} {\t group name} must be a alpha-numeric name (which may also 
contain underscores) which is unique 
within the scope of the thorn. A group name is compulsory.
\item{} {\t TYPE} designates the kind of variables held by the group.
The choices are {\t GF}, {\t ARRAY} or {\t SCALAR}. This field is 
optional, with the default variable type being {\t SCALAR}.
\item{} {\t DIM} defines the spatial dimension if the group is of type
{\t ARRAY} or {\t GF}, and can take the value {\t 1}, {\t 2}, or {\t 3}.
The default value is {\t DIM=3}.
\item{} The block following the group declaration line contains a list of
        variables contained in the group. All variables in a group have 
        the same data type, variable type and dimension. The list can be 
        separated by spaces, commas, or new lines. The variable names 
        must be unique within the scope of the thorn. 
        A variable can only be a member of 
        one group. The block must be delimited by brackets on new lines.
        If no block is given after a group declaration line, a 
        variable with the same name as the group is created.
\item{} An optional description of the group can be given on the last line, 
        at the moment this description is not used for any purpose.
\end{itemize} 

\section{param.ccl}
\label{sec:pa}

The parameter configuration file consists of a list of 
{\it parameter object specification items} (OSIs) giving the type and
range of the parameter separated by optional
{\it parameter data scoping items} (DSIs) which detail access to the
parameter.

\subsection{Parameter Data Scoping Items}
{\t
\begin{verbatim}
<access>:
\end{verbatim}
}
The key word {\t access} designates that all parameter object specification
items up to the next parameter data scoping item are in the same 
protection or scoping class. {\t access} can take the values:
\begin{itemize}
\item{} {\t public} --- all thorns have access to public parameters
\item{} {\t protected} --- other thorns can have access to these 
                           parameters if they specifically request
                           it in their own param.ccl
\item{} {\t private} --- only your thorn has access to private parameters
\item{} {\t friend} --- in this case an {\t implementation} name must follow
                        the colon. It declare that all the parameters in 
                        the following scoping block are protected variables
                        from the specified {\t implementation}.
\end{itemize}


\subsection{Parameter Object Specification Items}
\label{sec:paobspit}

{\t
\begin{verbatim}
[EXTENDS] <parameter type> <parameter name> ``<parameter description>'' 
{
  <PARAMETER VALUES>
} <default value>
\end{verbatim}
}
\begin{itemize}
\item{} Allowed {\t parameter types} are 
  \begin{itemize}
     \item{} {\t INTEGER} The specification of parameter values takes 
     the form of any number of comma-separated blocks of the form:
{\t
\begin{verbatim}
[<low-range>][:[<high-range>][:[<step>]]][::``<range description>'']
\end{verbatim}
}
Where an empty field, or a {\t *} in the place of {\t low-range} or
{\t high-range} indicates $-\infty$ and $\infty$ respectively. The 
default value for {\t step} is 1. A number by itself denotes that
this number is the only acceptable value.
     \item{} {\t REAL} The range specification is the same as integers,
             except that here, no {\t step} implies a continuum of values.
     \item{} {\t KEYWORD} Each entry in the list of acceptable values 
             for a keyword has the form 
{\t
\begin{verbatim}
``<keyword value>'', ``<keyword value>'' :: ``<description>''
\end{verbatim}
}
     \item{} {\t STRING}  No allowed values should be specified for a logical 
             parameter.
     \item{} {\t LOGICAL} No allowed values should be specified for a logical 
             parameter. The default value for a logical can be 
        \begin{itemize}
           \item{} True: {\t 1}, {\t yes}, {\t y}, {\t t}, {\t true}
           \item{} False: {\t 0}, {\t no}, {\t n}, {\t f}, {\t false}
        \end{itemize}
  \end{itemize}

\item{} The {\t parameter name} must be unique within the scope of the thorn.

\item{} A thorn can declare that it {\t EXTENDS} a parameter that it is a 
        friend of. This allows it to declare additional acceptable values.
        By default it is acceptable for two thorns to declare the same 
        value as acceptable. 
\end{itemize}


\section{schedule.ccl}
\label{sec:sc}

The schedule configuration files consists of 
\begin{itemize}
\item{} assignments to switch on storage and communications 
        for array variables at the start of program execution.
\item{} {\it schedule blocks} which schedule a subroutine from the
        thorn to be called at a specified time during program 
        executation. Statements within the schedule block can
        be used to switch on storage and communication for groups
        of variables during the duration that the subroutine is called.
\item{} Conditional statements
\end{itemize}

{\it assignments statements} have the form:
{\t
\begin{verbatim}
[STORAGE: <group>, <group>]
[COMMUNICATION: <group>, <group>]
\end{verbatim}
}

Each {\it schedule block} in the file {\t schedule.ccl} must have
the syntax:
{\t
\begin{verbatim}
schedule <function name> at <time>
{
  LANG: <language>
  [STORAGE: <group>,<group>]
  [COMMUNICATION: <group>,<group>]
  [TRIGGER: <grid function>,<grid function>]
} ``Description of function''
\end{verbatim}
}

Any schedule block or assignment statements can be optionally 
surrounded by conditional {\t if-elseif-else}
constructs using the parameter data base. These can be nested,
and have the general form:
{\t
\begin{verbatim}
if (CCTK_Equals(<parameter>,<string>)) 
{
  [<assignments>]
  [<schedule blocks>]
}
\end{verbatim}
}

Conditional constructs cannot be used inside of a schedule block.


\chapter{CCTK Parameters}
\label{sec:ccpa}

\chapter{Using GNATS}
GNATS is a freely redistributable set of tools for tracking bug
reports. It allows users to categorize their problem report and submit 
them to the GNATS. The bugtracker will assign appropriate maintainers
to analyze and solve the problem.
We are currently supporting a web based interface at {\tt
http://www.cactuscode.org} which lets you
interactively file a bug report. Here, we briefly
describe the main categories when creating a Cactus
problem report.
\begin{Lentry}
\item[{\bf Reporters email}] Your email address so we can get in
contact with you.
\item[{\bf Category}] there is currently a category for each of the
Cactus thorns and packages, also a category for the old Cactus3.x and
some general subjects like Web,etc. Select whatever category fits
best.
\item[{\bf Synopsis}] A brief and informative subject line.

\item[{\bf Confidential}] Unused, all PRs are public.

\item[{\bf Severity}] Pick one three levels.

\item[{\bf Class}] In the selected category, specify what kind of
problem you are dealing with.

\item[{\bf submitter ID}] Unused

\item[{\bf Originator}] Your name

\item[{\bf Release}] The Cactus release your are using.
%Type {\tt gmake version} FIXME

\item[{\bf Environment}] Very important: specify the environment:
e.g. {\tt uname -a}, MPI/non-MPI, etc.

\item[{\bf Description}] Describe your problem precisely, if you get a 
core dump, include the stack trace, give the minimal number of thorns, 
this problems occurs with.

\item[{\bf How-To-Repeat}] Tell us how to repeat the problem, if it is
software related.

\item[{\bf Fix}] If you can provide a fix, let us know.
\end{Lentry}

We also provide the customized {\tt send-pr} and {\tt send-pr.el} programs at
our website. These commands are compiled to submit Cactus problem
reports in your shell and from within emacs, respectively.


\label{sec:usgn}

\chapter{Using CVS}
\label{sec:uscv}
CVS is a version control system, which allows you to  keep
old versions of files (usually source code), keep a log of
who, when, and why changes occurred,  etc.
Unlike the simpler systems, CVS does not just operate on one file at a
time or one directory at a time,  but 
operates  on  hierarchical collections of directories consisting of
version controlled files.  CVS helps to  manage 
releases  and  to control the concurrent editing of source
files among multiple  authors. CVS can be obtained from 
{\tt http://www.cyclic.com}. 

A {\em CVS repository}located on a {\em server} may consist of an arbitrary
number of {\em modules}, which can be checked out independently. The Cactus
packages are organized as modules, their {\em CVS server} is {\tt cvs.cactuscode.org}.

\subsection{Essential CVS commands}

\begin{Lentry}
\item[{\bf cvs login}]
Logs into the repository. You will be prompted for a {\em
password}. This cvs command leaves a file {\tt .cvs-pass} in your
home directory. There is no need to login every time you issue a cvs 
command, as long as this file exists. For a Cactus checkout, you have
to log into the CVS server, using the cvs option {\bf -d} to specify CVSROOT:\\
{\tt cvs -d :pserver:cvs\_anon@cvs.cactuscode.org:/usr/users/cactus login}

Note, that there is no ``logout'' command: if you log in with
administrative rights from an arbitrary account, you should be aware
that the password file enables subsequent administrative logins from
that account. {\em Delete the file if necessary}.

\item[{\bf cvs checkout} {\em modules} \ldots]
This command  creates
your private copy of the source for {\em modules}. You can work
with this copy  without  interfering  with  others'
work.   At  least  one subdirectory level is always created: it does
not write all files into your current directory but creates a
directory. For Cactus, you need to include the {\bf -d} options to
specify the CVSROOT directory and the CVS server (see below). Once you 
have checked out the repository there is no need to include the {\bf
-d} option and its rather lengthy argument: the necessary information
is contained in the local {\tt CVS/} directories.

\item[{\bf cvs update}]
Execute  this  command  from  {\it within}  your  private
source  directory  when  you  wish  to  update your
copies of source  files  from  changes  that  other
developers have made to the source in the repository.
Merges are performed automatically when possible, a warning is issued
if manual  resolution  is  required  for conflicting  changes. If you
local copy is several version behind the actual repository copy, CVS
will {\em refetch} the whole file instead of applying multiple
patches.

\item[{\bf cvs add} {\tt file}]
Use this command to enroll new files in cvs records
of your working directory.  The files will be added
to the  repository  the  next  time  you  run  `cvs
commit'. 

\item[{\bf cvs commit} {\tt file}]
Use this command to add your local changes to the source to
the repository and thereby making it publically available to
checkouts and updates by other users. You cannot commit a
newly created file unless you have {\it added} it.

\item[{\bf cvs diff} {\tt file}]
Show differences between files in working directory
and  source repository, or between two revisions in
source repository.  (Does not change either repository or working
directory.) 

\item[{\bf cvs remove} {\tt file}]
Remove files from the source repository, pending  a {\tt cvs commit} on
the same files.  

\item[{\bf cvs status}] 
This command returns the current status of your local copy relative to 
the repository: e.g. it indicates local modifications and possible
updates.

\item[{\bf cvs import} {\tt repository tag1 tag2}]
Import adds an entire source distribution to the repository directory.
Use this command to add new packages to the Cactus4.0 repository. The
{\tt repository} argument gives a directory name (or a path to a
directory) and the CVS root directory for repositories; to obtain this 
directory on the CVS server, send a request to {\tt
cactus@cactuscode.org}. {\tt tag1} and {\tt tag2} are two tags (vendor 
 and release tags) that have to be supplied.

\end{Lentry}

\subsection{CVS Options}
The CVS command line can include the following:
\begin{itemize}
\item{\bf cvs options}, which apply to the overall cvs program
\item{\bf a cvs command}, which defines a particular task carried out by
cvs
\item{\bf command options} to specify certain working modes for the cvs
command.
\item{\bf command arguments} to specify which file to act on.
\end{itemize}

The options must be put {\em relative} to the {\it cvs command} as the
same option name can mean different things: {\em cvs options} go to the
{\em left} of the cvs command, {\em command options} go to the {\em right}
of the cvs command. Here is a list of essential {\em cvs options}

\begin{Lentry}

\item[{\bf -d} {\em cvs\_root\_directory}]
Use {\em cvs\_root\_directory} as the root directory  pathname  of
the  master source repository.  Overrides 
the setting of the  {\em CVSROOT}  environment  variable.
This value should be specified as an absolute pathname. 
In the Cactus4.0 checkout procedure, you specify the Cactus CVS server:\\
{\tt -d :pserver:cvs\_anon@cvs.cactuscode.org:/usr/users/cactus/}

\item[{\bf -z} {\em compression-level}]
When transferring files across the network use {\tt gzip}
with compression level  {\em compression-level}  to  compress  and
decompress  data as it is transferred. 
Requires the presence of the GNU gzip  program  in
the current search path at both ends of the link.

\item[{\bf -n}]
Do not change any file. Attempt to execute the {\em cvs command} but
only to issue reports. Does not remove, update, etc. any files. Very
effective for testing.

\item[{\bf -v}]
Displays version information of the installed CVS.

\item[{\bf -H} {\em cvs-command}]
Displays usage information about the specified cvs command. Without
cvs-command, a list of all available commands is returned. 
\end{Lentry}

Here is a list of essential {\em command options} with the
commands they are used with. They go after the cvs command. For a more 
complete list of all options, please refer to the manual page.

\begin{Lentry}

\item[{\bf -P}]
Prune  (remove)  directories  that  are empty after being updated, on
{\bf checkout}, or  {\bf update}.   Normally, an  empty  directory
(one that is void of revision controlled files) is  left  alone.
Specifying  {\bf -P} will cause these directories to be silently
removed from  your  checked-out  sources.   This  does  not  remove
the directory from the repository, only from your checked out copy.

\item[{\bf -m} {\em "Text"}]
Specify a logging message explaining changes, etc. on {\bf commit},
{\bf import}. If you do not specify a message, your default editor
is invoked to allow you to enter one.

\item[\bf -d]
Use this option with the {\bf update} command to create any
directories if they are missing from your local copy. This is normally 
the case if another user has added files or directories to the
repository. By default the {\bf update} command only acts on files in
your local copy. Note that omitting this option is a frequent cause of 
files missing during compilation.

\end{Lentry}

\subsection{CVS Examples}
We list some sample CVS commands to treat the most typical Cactus4.0
CVS situations. 
\begin{description}
\item\textbf{Logging into the server}\newline
{\tt cvs -d :pserver:cvs\_anon@cvs.cactuscode.org:/usr/users/cactus
login} \\ You will be asked for the password for user {\em cvs\_anon}, which is {\tt anon}.

\item\textbf{Checking out the code}\newline
{\tt cvs -d :pserver:cvs\_anon@cvs.cactuscode.org:/usr/users/cactus
checkout CCTK}\\
check out a CVS module named ``CCTK'', in this case it checks out the
Cactus Computational Toolkit. A directory {\tt ./CCTK} is created if  
it doesn't already exits. If you perform a checkout on an already
existing and locally modified copy of the module, CVS will try to merge the files
with your local copy. 

\item\textbf{Updating a file or directory}\newline
Assuming that you have a file {\tt ./foobar} in your checked out
copy, you may perform a \\
{\tt cvs status ./foobar}\\
to inform yourself about the necessary updates, etc. To update the
file issue \\
{\tt cvs update ./foobar}\\
If that was file was locally modified, CVS will try to merge the
changes.Manual merging might be necessary and will be indicated by a
CVS warning.

\item\textbf{Updating a directory}\newline
To recursively update the current directory and all subdirectories,
type\\
{\tt cvs update .}\\
To update a directory {\tt ./mysources}, type\\
{\tt cvs update ./path/to/mysources}

\item\textbf{Committing a changed file}\newline
To commit changes you have applied to your local copy, your file must be in
sync with the repository: your changes must be done to the
latest version, otherwise CVS will instruct you to perform an {\bf
update} first. To commit changes made to a file {\tt ./foobar}, type\\
{\tt cvs commit -m ``Reason for the change'' ./foobar}\\
You may specify several files to commit.

\item\textbf{Adding and committing a new file}\newline
Adding a new file to the repository is a two fold procedure you first
schedule the file for addition, then you commit it:\\
{\tt cvs add ./newfoo}\\
{\tt cvs commit -m ``new few message'' ./newfoo}

\item\textbf{Creating a new thorn}\newline
To add a new module to the Cactus repository, (e.g. an additional Cactus
package), you have to obtain a directory on the CVS server and a
login, which gives you administrative rights on your module. Please
contact the {\tt cactus@cactuscode.org} for details. The following
CVS steps apply as well to adding to a locally hosted CVS repository (except 
that the CVSROOT option has to be changed to your local host).
To perform the addition, type\\
{\tt cvs -d :pserver:}{\em your\_login}{\tt
@cvs.cactuscode.org:/usr/users/cactus }{\em package\_name} {\tt start V1}\\
where {\tt start} and {\tt V1} are the vendor and release tags, which
have to be supplied. 
\end{description}


\chapter{Using TAGS}
\label{sec:usta}
Finding your way around in the Cactus structure can be pretty
difficult to handle. To make life easier there is support for TAGS,
which lets you browse the code easily from within Emacs/XEmacs or vi. 
A TAGS database can be generated with by gmake:

\section{TAGS with emacs}
\label{sec:tawiem}

{\tt gmake TAGS} will create a database for a routine reference
table to be used within Emacs. This database can be accessed within
emacs if you add either of the following lines to your {\tt .emacs} file:\\
{\tt (setq tags-file-name "CCTK\_dir/TAGS")} XOR \\
{\tt (setq tag-table-alist '(("CCTK\_dir" . "CCTK\_dir/TAGS")))}\\
where {\tt CCTK\_dir } is your CCTK directory.\\

You can now easily navigate the CCTK and Toolkits by searching for
functions or ``tags'':
\begin{enumerate}
\item \textbf{ Alt.} will find a tag
\item \textbf{ Alt,} will find the next matching tag
\item \textbf{ Alt*} will go back to the last matched tag
\end{enumerate}
If you add the following lines to your {\tt .emacs} file, the
files found with TAGS will opened in {\em read-only} mode:
\begin{verbatim}
(defun find-tag-readonly (&rest a)
  (interactive)
  (call-interactively `find-tag a)
  (if (eq nil buffer-read-only) (setq buffer-read-only t))  )

(defun find-tag-readonly-next (&rest a)
  (interactive)
  (call-interactively `tags-loop-continue a)
  (if (eq nil buffer-read-only) (setq buffer-read-only t))  )

(global-set-key [(control meta \.)] 'find-tag-readonly)
(global-set-key [(control meta \,)] 'find-tag-readonly-next)
\end{verbatim}
The Key strokes to use when you want to browse in read-only mode are:
\begin{enumerate}
\item{CTRL Alt.} will find a tag and open the file in read-only mode
\item{CTRL Alt,} will find the next matching tag in read-only mode
\end{enumerate}

\section{TAGS with vi}
\label{sec:tawivi}

The commands available are highly dependent upon the version of {\bf vi}, but
the following is a selection of commands which may work.


\begin{enumerate}

\item \textbf{vi -t tag} 
Start vi and position the cursor at  the  file and line where `tag' is defined.

\item \textbf{Control-]} 
Find the tag under the cursor.

\item \textbf{:ta tag} 
Find a tag.

\item \textbf{:tnext} 
Find the next matching tag.

\item \textbf{:pop} 
Return to previous location before jump to tag.

\item \textbf{Control-T} 
Return to previous location before jump to tag (not widely implemented).

\end{enumerate}



\chapter{Thorn-Flesh Interface}
\label{sec:thflin}


%%%%%%%%%%%%%%%%%%%%%%%%%%%%%%%%%%%%%%%%%%%%%%%%%%%%%%%%%%%%%%%%%%%

\end{cactuspart}







